\documentclass[10pt,a4paper]{amsart}
\usepackage{amsmath,amssymb,amsthm,graphicx}
\usepackage{verbatim}
\newenvironment{reference}{\comment}{\endcomment}
%\newenvironment{reference}{}{}
\newenvironment{slogan}{\comment}{\endcomment}
\newenvironment{history}{\comment}{\endcomment}

% For commutative diagrams we use Xy-pic
\usepackage[all]{xy}

% We use 2cell for 2-commutative diagrams.
\xyoption{2cell}
% Theorem environments.
%
\theoremstyle{plain}
\newtheorem{theorem}[subsection]{Theorem}
\newtheorem{proposition}[subsection]{Proposition}
\newtheorem{lemma}[subsection]{Lemma}

\theoremstyle{definition}
\newtheorem{definition}[subsection]{Definition}
\newtheorem{example}[subsection]{Example}
\newtheorem{exercise}[subsection]{Exercise}
\newtheorem{situation}[subsection]{Situation}

\theoremstyle{remark}
\newtheorem{remark}[subsection]{Remark}
\newtheorem{remarks}[subsection]{Remarks}

\numberwithin{equation}{subsection}

% Macros
%
\def\lim{\mathop{\mathrm{lim}}\nolimits}
\def\colim{\mathop{\mathrm{colim}}\nolimits}
\def\Spec{\mathop{\mathrm{Spec}}}
\def\Hom{\mathop{\mathrm{Hom}}\nolimits}
\def\Ext{\mathop{\mathrm{Ext}}\nolimits}
\def\SheafHom{\mathop{\mathcal{H}\!\mathit{om}}\nolimits}
\def\SheafExt{\mathop{\mathcal{E}\!\mathit{xt}}\nolimits}
\def\Sch{\mathit{Sch}}
\def\Mor{\mathop{\mathrm{Mor}}\nolimits}
\def\Ob{\mathop{\mathrm{Ob}}\nolimits}
\def\Sh{\mathop{\mathit{Sh}}\nolimits}
\def\NL{\mathop{N\!L}\nolimits}
\def\CH{\mathop{\mathrm{CH}}\nolimits}
\def\proetale{{pro\text{-}\acute{e}tale}}
\def\etale{{\acute{e}tale}}
\def\QCoh{\mathit{QCoh}}
\def\Ker{\mathop{\mathrm{Ker}}}
\def\Im{\mathop{\mathrm{Im}}}
\def\Coker{\mathop{\mathrm{Coker}}}
\def\Coim{\mathop{\mathrm{Coim}}}

% Boxtimes
%
\DeclareMathSymbol{\boxtimes}{\mathbin}{AMSa}{"02}

%
% Macros for moduli stacks/spaces
%
\def\QCohstack{\mathcal{QC}\!\mathit{oh}}
\def\Cohstack{\mathcal{C}\!\mathit{oh}}
\def\Spacesstack{\mathcal{S}\!\mathit{paces}}
\def\Quotfunctor{\mathrm{Quot}}
\def\Hilbfunctor{\mathrm{Hilb}}
\def\Curvesstack{\mathcal{C}\!\mathit{urves}}
\def\Polarizedstack{\mathcal{P}\!\mathit{olarized}}
\def\Complexesstack{\mathcal{C}\!\mathit{omplexes}}
% \Pic is the operator that assigns to X its picard group, usage \Pic(X)
% \Picardstack_{X/B} denotes the Picard stack of X over B
% \Picardfunctor_{X/B} denotes the Picard functor of X over B
\def\Pic{\mathop{\mathrm{Pic}}\nolimits}
\def\Picardstack{\mathcal{P}\!\mathit{ic}}
\def\Picardfunctor{\mathrm{Pic}}
\def\Deformationcategory{\mathcal{D}\!\mathit{ef}}
\usepackage{cleveref}
\crefname{originallemma}{original lemma}{original lemmas}
\Crefname{originallemma}{Original lemma}{Original lemmas}
\makeatletter
\newlabel{topology-lemma-closed-union-connected-components}{{04PL}{source}{}{sourcepointer-1}{}}
\newlabel{topology-lemma-closed-union-connected-components@cref}{{[originallemma][0][]04PL}{[1][0][]source}{}{}{}}
\makeatother
\title{Realizing the inverse image of a closed subset of the components as an ind-Zariski quotient}
\author{The Stacks Project Authors}
\date{stacks-project:proetale.tex:local-snapshot-20261012; GNU FDL 1.2 or later}
\begin{document}
\maketitle
\noindent\textit{Editorial note.} All delivered mathematical text is the author's own, copied from the pinned Stacks Project chapter file proetale.tex; no statement or proof step was rewritten, reordered or removed.\par\medskip
\noindent\textit{Editorial note.} The counted proof is the author's complete argument: the presentation of the inverse image of the closed subset as an intersection of open and closed subsets, the description of each piece as the spectrum of a localization at an idempotent, and the colimit construction with the resulting homeomorphism.\par\medskip
\noindent\textit{Editorial note.} The label lemma-construct names the counted statement, and the proof environment immediately following it in the source is delivered verbatim as its proof.\par\medskip
\noindent\textit{Editorial note.} The counted statement is carried by the author's own lemma environment in the chapter source.\par\medskip
\noindent\textit{Editorial note.} Pre-open gate: python operations/agent-authoring/tools/precard\_gate.py --source proetale.tex --label lemma-construct --start 1080 --end 1095 --local-id 28009 returned RESULT=OK (L28009 free, occupancy FREE, public interval clear). Candidate taken from sources/stacks-proetale/candidates-halfB.json.\par\medskip
\noindent\textit{Editorial note.} The statement and proof are copied from the pinned chapter file; the cited inputs are kept as external pointers with their canonical Stacks tags.\par\medskip
\noindent\textit{Allowed standard prerequisites (editorial).}\begin{itemize}
\item Closed subsets of the space of connected components as intersections of open and closed subsets. Topology, the lemma describing a closed subset of the space of connected components as the intersection of the open and closed subsets containing it. Supplies the presentation of the inverse image Z of the closed subset T as an intersection of open and closed subsets Z\_alpha.
\item Open and closed subsets of an affine scheme as spectra of localizations at idempotents. Algebra, the description of the open and closed subsets of the spectrum of a ring as the spectra of localizations at idempotents, giving local isomorphisms. Supplies the rings A\_alpha with A -> A\_alpha a local isomorphism and Spec(A\_alpha) = Z\_alpha.
\end{itemize}\medskip
\ifdefined\hypertarget\hypertarget{sourcepointer-1}{}\fi
\noindent\textit{Source reference.} Original lemma 04PL: sources/stacks-chow/tags.txt, source lines 6119--6119. Prerequisite pointer invoked at proof line(s) 1091; it names a standard fact of a sibling chapter that the counted argument uses, and every step of the counted proof built on it is retained verbatim in the delivered body.\par\medskip
\begin{lemma}
\label{lemma-construct}
Let $A$ be a ring. Let $X = \Spec(A)$. Let $T \subset \pi_0(X)$ be a
closed subset. There exists a surjective ind-Zariski ring map $A \to B$
such that $\Spec(B) \to \Spec(A)$ induces a homeomorphism of $\Spec(B)$
with the inverse image of $T$ in $X$.
\end{lemma}

\begin{proof}
Let $Z \subset X$ be the inverse image of $T$. Then $Z$ is the intersection
$Z = \bigcap Z_\alpha$ of the open and closed subsets of $X$ containing $Z$,
see Topology, Lemma \ref{topology-lemma-closed-union-connected-components}.
For each $\alpha$ we have $Z_\alpha = \Spec(A_\alpha)$ where
$A \to A_\alpha$ is a local isomorphism (a localization at an idempotent).
Setting $B = \colim A_\alpha$ proves the lemma.
\end{proof}

\end{document}
